% =====================================================================
%  CHAPTER 5: 公共营养
% =====================================================================
\chapter{公共营养}
\chapterintro{
\keyword{掌握}DRIs七项指标体系（EAR/RNI/AI/UL/AMDR/PI-NCD/SPL）的定义、制定逻辑和应用；掌握膳食模式与膳食指南（中国居民膳食指南2022八条核心准则与膳食宝塔2022构成及应用）；掌握营养调查膳食调查方法及膳食营养评价方法和依据；掌握BMI分级标准及NRV概念与应用。\keyword{熟悉}营养调查与评价（调查目的、内容、综合营养状况评价方法、营养监测概念及指标）；熟悉公共营养改善（营养改善措施、营养教育、营养干预的方法和应用）；熟悉食品营养标签定义内容及管理规范；熟悉功能食品概念分类及管理；熟悉食谱编制基本原则方法；熟悉全国营养调查概况。\keyword{了解}公共营养的定义与地位；了解我国营养相关政策与法规。
}

% =====================================================================
\section{概述}
% =====================================================================

\begin{defbox}
\textbf{公共营养的定义：}从医学状况+人群社会视角的群体概念。定义范围涵盖：病人、老人、慢性病患者、营养不良患者。人群分类：高危人群、亚健康人群、一般人群和特殊人群。

\textbf{合理营养与平衡膳食：}
\begin{itemize}[leftmargin=*]
    \item \textbf{合理营养（rational nutrition）：}指通过膳食提供给机体种类齐全、数量充足、比例适宜的能量和各种营养素，并与机体需要相平衡，从而达到合理营养、促进健康、预防疾病的膳食。又称为平衡膳食（balanced diet）。
    \item \textbf{平衡膳食的要求：}1. 提供种类齐全、数量充足、比例适宜的营养素；2. 保证食物安全；3. 科学烹调加工；4. 合理的进餐制度和良好的饮食习惯。
\end{itemize}

\textbf{营养不良（malnutrition）：}指由于一种或一种以上营养素的缺乏或过剩所造成的机体健康异常或疾病状态。营养不良包括两种表现：营养缺乏（nutrition deficiency）和营养过剩（nutrition excess）。
\end{defbox}

\begin{defbox}
\textbf{合理膳食与公共营养的对应关系：}
\begin{itemize}[leftmargin=*]
    \item 营养素 $\rightarrow$ 营养标准（DRIs）
    \item 食物 $\rightarrow$ 膳食指南/食物指导
    \item 膳食 $\rightarrow$ 膳食模式
\end{itemize}

\textbf{公共营养的地位与作用：}
\begin{enumerate}[leftmargin=*]
    \item 第一，促进全人群的营养与健康。
    \item 第二，增强国民体质，增强综合国力。
    \item 第三，为国家经济社会发展和民族复兴提供建议，事关国家发展的战略性问题。
\end{enumerate}
\end{defbox}

\begin{defbox}
\textbf{公共营养的工作内容：}
\begin{enumerate}[leftmargin=*]
    \item 营养标准的制定（DRIs）
    \item 居民膳食指南的制定
    \item 各类人员营养培训与指导
    \item 营养调查与评价
    \item 营养监测
    \item 营养干预
    \item 食物营养规划与营养立法
    \item 其他营养公众指导（以食物为研究对象、食物营养政策和法规）
\end{enumerate}

\textbf{公共营养的管理机构：}
\begin{itemize}[leftmargin=*]
    \item 全国营养工作的组织协调：中国营养学会
    \item 制定营养素标准：国家卫生健康委营养标准委员会
    \item 制定营养相关政策及法规：食品安全法、营养改善工作管理办法、全国营养调查工作规范、食品营养标签管理规范、《2025-2030年中国食物与营养发展纲要》、2017年7月《国民营养计划（2017-2030）》
\end{itemize}
\end{defbox}

% =====================================================================
\section{膳食营养素参考摄入量（DRIs）}\difficulty
% =====================================================================

% [补充自往届笔记] - DRIs制定基础说明
每个人每日从膳食中获取各种营养素及能量。对某种营养素的需要量因年龄、性别、生理状况而不同（如儿童、孕妇、老人、哺乳期妇女）。人们所需要的各种营养素都从每日膳食中获得。科学安排每日膳食才能满足人体的营养素需要。某种营养素长期摄入不足或摄入过多都会产生危害。因此针对各年龄、性别和劳动强度状态人群，制定了膳食营养素参考摄入量（dietary reference intakes, DRIs），用于指导膳食设计和评价膳食计划。

\begin{defbox}
\textbf{DRIs（Dietary Reference Intakes）：}一组每日平均膳食营养素摄入量的参考值。是评价膳食营养素供给量能否满足人体需要及是否存在过量摄入风险的一组参考值。

\textbf{从RDAs到DRIs的发展：}
\begin{itemize}[leftmargin=*]
    \item RDAs：营养目标——预防营养缺乏病
    \item DRIs：预防慢性病——帮助个体合理安排膳食、获得良好健康状态的新概念。用途：健康个体的膳食计划、膳食评估；并用于制定政策计划。
\end{itemize}
\end{defbox}

\begin{keybox}
\textbf{DRIs七项指标体系（DRIs 2013版）：}

\begin{table}[H]
\centering
\begin{tabularx}{\textwidth}{lXl}
\toprule
\textbf{指标} & \textbf{定义} & \textbf{制定目标} \\
\midrule
\imp{平均需要量（EAR）} & 某一特定性别、年龄及生理状况群体中对某营养素需要量的平均值。摄入量达到EAR水平时可以满足群体中50\%个体的需要，不能满足另外50\%个体对该营养素的需要。 & 预防缺乏 \\
\hline
\imp{推荐摄入量（RNI）} & 相当于RDA，可以满足某一特定性别、年龄及生理状况群体中绝大多数（97\%-98\%）个体需要量的摄入水平。长期摄入RNI水平，可以满足个体对该营养素的需要，维持组织中有适当的营养素储备和机体健康。RNI是以EAR为基础制定的。用途：作为个体每日摄入该营养素的目标值。 & 预防缺乏 \\
\hline
\imp{适宜摄入量（AI）} & 当某种营养素个体需要量的资料不足，没有办法计算出EAR，因而不能求得RNI时，可设定适宜摄入量来代替RNI。AI是通过观察或实验获得的健康人群某种营养素的摄入量，是该营养素的一个安全摄入水平，而不是通过研究营养素的个体需要量得到的。用途：作为个体每日营养素摄入的目标值。 & 预防缺乏 \\
\hline
\imp{可耐受最高摄入量（UL）} & 平均每日可以摄入营养素的最高限量。摄入量达到UL对一般人群中的几乎所有个体都不至于损害健康。在规划膳食时应使营养素摄入量低于UL，以避免造成不良反应的潜在风险。 & 防止过量 \\
\hline
\imp{宏量营养素可接受范围（AMDR）} & 指碳水化合物、脂肪和蛋白质的理想摄入量范围，以占能量摄入量的百分比(\%E)表示。该范围可提供人体对这些必需营养素的需要，并且有利于降低慢性病的发生危险。其下限为预防营养缺乏，上限为降低慢性非传染性疾病风险。如果一个个体的摄入量高于或低于推荐的范围，可能增加罹患慢性病的风险，或者导致必需营养素缺乏的可能性增加。 & 预防慢性病 \\
\hline
\imp{预防非传染性慢性病的建议摄入量（PI-NCD）} & 膳食营养素摄入过高导致的慢性病（肥胖、糖尿病、高血压、血脂异常、脑卒中、冠心病以及某些癌症等）。PI-NCD以非传染性慢性病的一级预防为目标，提出的必需营养素的每日摄入量。当NCD易感人群某些营养素的摄入量接近或达到PI值时，可以降低发生NCD的风险。某些营养素PI可能高于RNI或AI（如维生素C、钾）；另一些营养素可能低于AI（如钠）。 & 预防慢性病 \\
\hline
\imp{特定建议值（SPL）} & 研究证明某些食物成分具有改善生理功能、预防营养相关慢性病的生物学作用，其中多数为植物化学物。SPL以降低人群中膳食相关NCD风险为目标，提出其他食物成分的每日摄入水平，当该成分的摄入量达到SPL时，可能有利于降低疾病的发生风险或改善健康状况。 & 预防慢性病 \\
\bottomrule
\end{tabularx}
\end{table}
\end{keybox}

\begin{keybox}
\textbf{RNI的制定逻辑：}
RNI是以EAR为基础制定的。如果人群需要量呈正态分布，已知EAR的标准差，则RNI = EAR + 2SD表示。如果营养素需要量资料不符合正态分布，不能计算SD时，一般采用EAR乘以变异系数（10\%），即RNI = 1.2 × EAR。

\textbf{RNI设定注意事项：}RNI是根据某一特定人群中体重在正常范围内的个体的需要量设定的。对于个别身高、体重超出参考范围较多的个体，可能需要按每千克体重的需要量调整其RNI。
\end{keybox}

\begin{keybox}
\textbf{AI与RNI的异同：}
\begin{itemize}[leftmargin=*]
    \item \textbf{相同点：}都可以作为群体中个体营养素摄入量的目标值，可以满足群体中几乎所有个体的需要。
    \item \textbf{不同点：}AI的准确性不及RNI，主要用于个体营养素摄入目标；群体中营养素平均摄入量达到或超过AI，则该人群中摄入不足者的危险性很小。
\end{itemize}
\end{keybox}

\begin{keybox}
\textbf{UL的制定注意事项：}
\begin{itemize}[leftmargin=*]
    \item 随着营养素强化食品和膳食补充剂的广泛使用，需要制定ULs以指导安全消费。
    \item 如果某营养素的毒副作用与摄入总量有关，则该营养素的UL值根据食物、饮水及补充剂提供的总量来制定。
    \item 如果毒副作用仅与强化食物和补充剂有关，则UL根据这些来源的总摄入量来制定。
    \item 对于某些营养素而言，没有足够资料来制定其UL，但未定UL并不意味着过多摄入没有潜在的危害。
\end{itemize}
\end{keybox}

\begin{tipbox}
\textbf{DRIs的制定与修订程序：}
人群健康与营养状况调查和监测 $\rightarrow$ 制定DRIs $\rightarrow$ 膳食指导 $\rightarrow$ 食物指导 $\rightarrow$ 预包装食品营养标签 $\rightarrow$ 人群营养改善和健康状况监测反馈 $\rightarrow$ 营养科学与有关科学研究（实验室、临床观察、流行病学调查）$\rightarrow$ 修订DRIs。注意：每4-5年修订一次。

\textbf{DRIs的用途与层次：}
\begin{itemize}[leftmargin=*]
    \item \textbf{个体用途：}膳食规划、制定食谱、食品配方研发、食物保障/供给计划。
    \item \textbf{群体用途：}膳食评价（膳食评估和改进、营养咨询、营养风险）。
    \item \textbf{安全作用：}作为标准/依据，衔接个体膳食规划和群体膳食评价。
\end{itemize}

\textbf{评价逻辑：}个体营养素需要量 VS 膳食营养素供给量/摄入量 $\leftrightarrow$ 膳食营养素参考摄入量（DRIs）
\end{tipbox}

\begin{defbox}
\textbf{能量需要量（EER, Estimated Energy Requirement）：}
EER是指能长期保持良好健康状态、维持体重和机体构成以及体力活动水平的个体及人群，达到能量平衡时所需要的膳食能量摄入量。

群体能量推荐摄入量直接等同于该群体的能量EAR，而不同于蛋白质等其他营养素采用EAR加上2个标准差的方法。

EER的制定需考虑年龄、性别、体重、身高和体力活动的不同：
\begin{itemize}[leftmargin=*]
    \item 成人EER：一定年龄、性别、体重、身高和体力活动的健康人群，维持能量平衡所需的膳食能量。
    \item 儿童EER：一定年龄、性别3岁以上儿童正常体重、身高的个体，维持能量平衡和正常生长所需膳食能量。
    \item 孕妇EER：在成人EER基础上，再加上胎儿组织生长所需的膳食能量。
    \item 乳母EER：在成人EER基础上，再加上乳汁分泌所需的膳食能量。
\end{itemize}
\end{defbox}

% =====================================================================
\section{膳食模式与膳食指南}
% =====================================================================

\subsection{膳食结构}

\begin{defbox}
\textbf{膳食结构（Dietary Structure）的定义：}
膳食中各类食物的种类、数量及其在膳食中所占的比例。影响因素：历史时期、民族/地区、社会阶层。特性：可变性，可预见性和相对稳定性及代表性。核心构成：粮食、动物、其他。

\textbf{膳食结构研究的意义：}
\begin{itemize}[leftmargin=*]
    \item 研究一个国家或一个地区人群膳食特点和营养状况的基础。
    \item 比较差别和变化，了解发展、饮食文化、习惯的变化规律、可能的健康风险。
    \item 为制定相应的膳食指导和营养指导提供科学依据。
\end{itemize}
\end{defbox}

\begin{defbox}
\textbf{膳食模式 $\approx$ 膳食结构：}
将日常膳食中成分组合而构成的独特食物和营养素摄入的方式，即膳食的总体构成。

\textbf{素食（VU膳食）定义：}指膳食中不含畜肉、禽肉、鱼类及其制品，且不含动物性制品和动物油脂的饮食习惯。其膳食构成不同于一般膳食，能量密度较低，富含DF（膳食纤维）、镁、叶酸、维生素C、E、n-6 PUFA（多不饱和脂肪酸）、植物化学物和较低的胆固醇、SFA（饱和脂肪酸）。
\end{defbox}

\begin{keybox}
\textbf{膳食结构类型及与健康的关系：}

\begin{enumerate}[leftmargin=*]
    \item \textbf{动物性食物为主的膳食结构（西方模式/经济发达国家模式）}——欧美经济发达国家。
    \begin{itemize}
        \item 营养过剩型，"三高"一低（高能量、高脂肪、高蛋白、低膳食纤维）。
        \item 健康风险：肥胖、心脑血管疾病等代谢性疾病、结肠癌、直肠癌等。
    \end{itemize}

    \item \textbf{植物性食物为主的膳食结构（东方模式/发展中国家模式）}——亚洲、非洲部分国家/地区。
    \begin{itemize}
        \item 营养缺乏型，"一高三低"（低蛋白质、低脂肪、植物性食物为主）。
        \item 健康风险：感染性疾病增加、营养缺乏症与发达国家类似的慢性病同时并存。
    \end{itemize}

    \item \textbf{动植物食物平衡的膳食结构}——日本为代表，集合东西方膳食特点。
    \begin{itemize}
        \item 动植物消费均衡、水产品食物量大，三大宏量营养素供能比合理。
        \item 近40年也在改变中，水产品消费量依然在增加，动物性食物消费量增加。
    \end{itemize}

    \item \textbf{地中海膳食结构}——希腊、西班牙、法国和意大利南部等。
    \begin{itemize}
        \item 橄榄油为主要脂肪来源（不饱和脂肪高），大量蔬菜水果、全谷物、豆类，适量鱼类和红酒，饱和脂肪低。
        \item 大部分成年人有饮用红酒的习惯。
    \end{itemize}
\end{enumerate}
\end{keybox}

\begin{keybox}
\textbf{我国膳食结构：}
\begin{itemize}[leftmargin=*]
    \item 地域偏差：营养缺乏与某些营养素过剩并存。
    \begin{enumerate}[leftmargin=*]
        \item 少部分地区仍以贫困型膳食结构为主。
        \item 沿海城市等地区已经呈现富裕型膳食结构。
        \item 大部分城市和地区正朝富裕型膳食结构的方向转变。
        \item 营养缺乏与营养过剩双重负担并存。
    \end{enumerate}
    \item \textbf{营养失衡表现：}营养缺乏/不足（缺乏症）；营养过剩/过量（肥胖、慢性病）。
\end{itemize}
\end{keybox}

\subsection{膳食指南}

\begin{defbox}
\textbf{膳食指南（Dietary Guidelines）的定义：}
根据营养学原理，制定的一组以食物为基础的膳食行为建议，用于指导大众合理选择与搭配食物；将营养素优化与化学行为相结合，倡导合理膳食维持健康，预防膳食相关疾病。% [补充自往届笔记] DG是根据营养科学原理，结合当地食物供应和大众需求，针对实际存在的营养问题，由营养健康权威机构制定的食物选择和身体活动的指导性意见。

\textbf{膳食指南对群体的核心要求（八大特征）：}
\begin{enumerate}[leftmargin=*]
    \item 群体性：针对市民/大众，解决群体性问题，具有针对性
    \item 科学性：倡导平衡膳食，改善营养状况，预防膳食相关疾病
    \item 普及性：采用通俗易懂的大众语言，强调可操作性
    \item 有效性：核心膳食建议经科学验证有效
    \item 通俗性：使用通俗易懂的大众语言
    \item 权威性：由权威机构制定并指导
    \item 保障性：确保食物供应
    \item 预防性
\end{enumerate}
\end{defbox}

\begin{keybox}
\textbf{中国居民膳食指南（2022）——8条核心准则（针对2岁以上健康人群）：}

\begin{enumerate}[leftmargin=*]
    \item \textbf{准则一：食物多样，合理搭配。}谷类为主，平均每天摄入12种以上食物，每周25种以上。
    \item \textbf{准则二：吃动平衡，健康体重。}
    \item \textbf{准则三：多吃蔬果、奶类、全谷、大豆。}奶及奶制品300-500 ml（液态奶）。
    \item \textbf{准则四：适量吃鱼、禽、蛋、瘦肉。}平均每天120-200 g，每周最好吃鱼2次。
    \item \textbf{准则五：少盐少油，控糖限酒。}每天食盐不超过5 g，糖不超过50 g（最好25 g以下）。
    \item \textbf{准则六：规律进餐，足量饮水。}女性1500 ml，男性1700 ml（8杯），提倡喝白开水或茶水。
    \item \textbf{准则七：会烹会选，会看标签。}
    \item \textbf{准则八：公筷分餐，杜绝浪费。}
\end{enumerate}
\end{keybox}

\begin{keybox}
\textbf{中国居民平衡膳食宝塔（2022）：}

\begin{table}[H]
\centering
\begin{tabularx}{\textwidth}{clX}
\toprule
\textbf{层级} & \textbf{食物类别} & \textbf{每日推荐量} \\
\midrule
第一层 & 谷薯类 & 250-400 g/d（全谷物和杂豆50-150 g，薯类50-100 g） \\
\midrule
第二层 & 蔬菜类 & 300-500 g/d（深色蔬菜占1/2） \\
       & 水果类 & 200-350 g/d \\
\midrule
第三层 & 鱼、禽、肉、蛋 & 120-200 g/d（鱼40-75 g，畜禽肉40-75 g，蛋类1个约50 g） \\
\midrule
第四层 & 奶及奶制品 & 300-500 g/d \\
       & 大豆及坚果 & 25-35 g/d \\
\midrule
第五层 & 盐 & $<$ 5 g/d \\
       & 油 & 25-30 g/d \\
\midrule
\multicolumn{2}{l}{\textbf{其他推荐：}} & 饮水1500-1700 ml/d，活动6000步/d \\
\bottomrule
\end{tabularx}
\end{table}
\end{keybox}

\begin{keybox}
\textbf{特定人群膳食指南主要原则（熟悉）：}

\begin{enumerate}[leftmargin=*]
    \item \textbf{备孕妇女：}调整孕前体重至适宜水平；常吃含铁丰富食物，选用碘盐；孕前3个月开始补充叶酸；禁烟酒，保持健康生活方式。
    \item \textbf{孕期妇女：}补充叶酸，常吃含铁丰富食物，选用碘盐；孕吐严重者少量多餐，保证必要量碳水化合物摄入；孕中晚期适量增加奶、鱼、禽、蛋、瘦肉摄入；适量体力活动，维持孕期适宜体重。
    \item \textbf{哺乳期妇女：}增加富含优质蛋白质及维生素A的动物性食物和海产品摄入；适量增加奶类，多喝汤水；愉悦心情，充足睡眠，坚持哺乳。
    \item \textbf{6月龄内婴儿：}坚持纯母乳喂养；生后1小时内开奶，重视尽早吸吮；回应式喂养；补充维生素D，无需补钙。
    \item \textbf{7-24月龄婴幼儿：}继续母乳喂养，满6月龄添加辅食，从富铁泥糊状食物开始；引入多样化食物，重视动物性食物；保持食物原味，不加盐糖；回应式喂养。
    \item \textbf{学龄前儿童：}规律就餐，自主进食；每天饮奶，足量饮水；合理烹调，少调料少油炸；保证户外活动时间。
    \item \textbf{学龄儿童：}主动参与食物选择与制作；吃好早餐，合理选择零食；天天喝奶，足量饮水，不喝含糖饮料，禁止饮酒；多户外活动。
    \item \textbf{老年人：}食物品种丰富，动物性食物充足，常吃大豆制品；鼓励共同进餐；积极户外活动，延缓肌肉衰减；保持适宜体重，定期体检。
    \item \textbf{素食人群：}谷类为主，食物多样；增加大豆及其制品的摄入；常吃坚果、海藻和菌菇类；蔬菜水果充足；合理选择烹调油。
\end{enumerate}
\end{keybox}

% =====================================================================
\section{营养调查与评价}
% =====================================================================

% [补充自往届笔记]
营养调查（Nutrition Survey）是通过膳食调查、体格测量、营养水平生化检验和营养缺乏症临床检查等方法，全面了解个体或群体的营养状况的方法。营养状况综合评价包括四个层面：营养调查、体格测量、生化检验和临床检查。营养调查是对大规模人群膳食摄入量、临床检查、生化检测信息的综合解释。人体状况受各种营养素摄入量、需要量的综合影响，营养状况评价可以从整体上确定相关人群的健康状况。

\begin{defbox}
\textbf{营养调查与相关概念区分：}
\begin{itemize}[leftmargin=*]
    \item \textbf{营养调查（Nutrition Survey）：}指运用各种手段准确了解某一人群特定时间的各种营养指标水平，并判断其营养和健康状况。涉及一次性大规模横断面调查。
    \item \textbf{营养监测（Nutrition Surveillance）：}通过长期纵向动态监测目标人群营养状态。
    \item \textbf{营养筛查（Nutrition Screening）：}通过筛查方法快速识别群体中已患营养不良的个体。
\end{itemize}
\end{defbox}

\subsection{营养调查的目的与内容}

\begin{keybox}
\textbf{营养调查的目的：}
\begin{enumerate}[leftmargin=*]
    \item 了解不同地区、年龄、性别人群的能量和营养素摄入现状。
    \item 了解与能量和营养素摄入不足和过剩有关的营养问题的分布和严重程度。
    \item 探索营养相关疾病的病因和影响因素。
    \item 了解人群膳食结构变迁及变化趋势。
    \item 为某些营养相关国家性科学研究提供基础营养和健康数据。
    \item 为国家或地方制定食物营养政策、法规、标准、干预措施以及发展规划提供科学依据。
\end{enumerate}
\end{keybox}

% [补充自往届笔记]
为了解国民的营养和健康状况，许多国家定期有计划地开展国民营养调查工作。为研究不同经济发展时期人群的膳食营养和健康状况的变化提供基础资料。为食物生产加工、群众食物消费和制定干预措施提供依据。

\subsection{膳食调查}

\begin{defbox}
\textbf{膳食调查的定义：}
收集调查对象在一定时期内食物消费数据，通过食物消费信息了解其膳食摄入情况（能量和各种营养素的摄入量、食物来源、供能比例、加工方式和进餐制度等），并评价其能量和营养素摄入量满足机体需要的程度。

\textbf{膳食调查方法的分类维度：}
\begin{itemize}[leftmargin=*]
    \item 称重法、记账法、询问法、食物频率法、化学分析法。
    \item \textbf{方法选择依据（5W1H）：}
    \begin{itemize}
        \item Who——研究对象（个体/群体）
        \item What——研究关注点（食物/营养素）
        \item When——关注当前/通常膳食模式/长期（一天/一周/一年）
        \item Where——在家/在外食用
        \item Why——研究目的
    \end{itemize}
    \item 依据：研究目的、研究人群大小、精确度要求、经费及研究时间。调查方法各有其优缺点，确定后统一使用，不能混合使用。
\end{itemize}
\end{defbox}

\begin{defbox}
\textbf{膳食调查方法的基本要求：}
\begin{itemize}[leftmargin=*]
    \item 观察单位：个体、家庭、群体
    \item 资料收集方式：记录法、询问法（面访、电话、自动应答、电子化技术、双份法）
    \item 研究时限：既往的、目前的
    \item 食物量估计：重量、体积、尺寸估计，图片和模型或工具
    \item 食物转换为营养素方法：营养素数据库（食物成分表）、直接的化学分析
\end{itemize}
\end{defbox}

\subsection{膳食调查的主要方法}

\begin{tipbox}
\textbf{营养调查方法分类总览（了解）：}

\textbf{一、膳食调查方法：}
\begin{itemize}[leftmargin=*]
    \item \textbf{个人：}24小时回顾法、食物频率法、食物称重记录法。
    \item \textbf{家庭：}家庭食物量登记和记账法、家庭食物称重法。
    \item \textbf{群体：}食物供应调查、食物平衡表、市场供应调查。
\end{itemize}
\end{tipbox}

\begin{defbox}
\textbf{膳食调查五大方法详解：}

\textbf{（一）24小时膳食回顾法（24h History Recalls）}

又称询问法，是获得个体食物摄入量最常用的方法。要求每个调查对象回顾和描述在调查时刻以前24h内摄入的所有食物（包括饮料）的种类和数量。一般选择3天连续调查方法，每天入户回顾24小时摄入情况，连续进行3天（含2个工作日+1个周末）。

\begin{itemize}[leftmargin=*]
    \item \textbf{优点：}不改变被调查者的饮食习惯；食物量能估计比较准确；调查时间短；一年中多次回顾可提供食物和营养素通常摄入量的估计值；应答者负担轻，应答率较高；操作简单、费用低。
    \item \textbf{缺点：}食物份额量大小不易准确判断；对某些人群可能有难度（如80岁以上老年人），可能不适用。
    \item \textbf{应用：}适用于人群群体的膳食调查和评估，以及个体患者的调查（如病人、散居儿童、老人、咨询门诊等）。
\end{itemize}

\textbf{（二）食物频率问卷法（Food Frequency Questionnaires, FFQ）}

是估计被调查者在规定的一段时期内摄入某些食物频率的方法，用于评估长期膳食营养状况。调查期可从几天、1周、1个月到3个月或1年以上。食物频率问卷由调查对象的基本信息和食物频率两部分组成。

\begin{itemize}[leftmargin=*]
    \item \textbf{FFQ分类：}
    \begin{itemize}
        \item \textbf{定性FFQ：}通常只得到每种食物在特定时期内（如过去1个月）摄入的次数，不收集食物份额量大小资料。
        \item \textbf{定量FFQ：}通常要求受访者提供所摄入食物的平均数量，一般借助于测量辅助物。
        \item \textbf{半定量FFQ：}通常要求研究对象提供标准（或平均）食物份额量大小的种类，通常提供选项选择。
    \end{itemize}
    \item \textbf{优点：}能得到通常膳食摄入量；不需要训练有素的调查员；调查方法简单、费用低；习惯性膳食模式不受影响；应答者负担轻，应答率高；能获得通常膳食摄入的数据（某些食物和营养素的量），更有助于研究膳食与疾病之间的关系。
    \item \textbf{缺点：}需要对过去膳食模式的记忆，较长时期可能不准确；当前的膳食模式可能影响对过去膳食的回顾，导致偏倚；食物份额量标准大小的估计不准；食物摄入量的估计可能不准确；是一种复杂的方法，对非营养专业人士而言，自行管理的调查表错误率高；应答者负担取决于调查食物种类的数量、复杂性和频次；通常不能测量食物摄入量的每日变异，难以确定个别调查结果是否准确。
    \item \textbf{说明：}24h回顾法和FFQ联合使用可以较全面地反映人群膳食摄入的结果，弥补询问调查法的局限。在膳食与健康关系的流行病学研究中应用最广泛。
\end{itemize}

\textbf{（三）称重法（Weighed Food Records）}

指运用各种称量工具对某一饮食单位（个人或家庭）所消耗的各种食物的生重、熟重及餐后剩余量进行称重，了解食物消耗情况，计算出每人每日各种营养素的平均摄入量。调查时间一般为3-7天。

\begin{itemize}[leftmargin=*]
    \item \textbf{优点：}能测定食物份额大小或称重；能准确反映被调查对象的食物摄入量，可调查每日每餐膳食的变动情况，能准确计算每人每日各种营养素的摄入量，可以作为膳食调查的"金标准"用于验证其他方法的准确性。
    \item \textbf{缺点：}需要较大人力物力，对调查员要求高，对应答者文化水平有一定要求，额外的负担可能造成应答率下降，不能保证样本的代表性。
    \item \textbf{应用：}适用于家庭、个人及特殊工作人员的膳食调查，不太适用于大规模普查工作（如流行病学调查）。
\end{itemize}

\textbf{（四）记账法（Account Method）}

通过记录一定时期内某一饮食单位（机关、学校、部队、单位食堂等）的食物消耗总量和进餐人数，来计算每人每日各种营养素的平均摄入量。可调查较长时期膳食（一个月或更长）。

\begin{itemize}[leftmargin=*]
    \item \textbf{优点：}调查目的准确，每餐用餐人数统计较确实的情况下，使调查结果较准确；可调查较长时期的膳食。
    \item \textbf{缺点：}调查结果的准确性取决于账目记录和人数统计的准确度；不能得出个体间营养摄入差异，不能用于分析个体膳食摄入状况。
    \item \textbf{应用：}适用于有较详细账目的集体单位和家庭。
\end{itemize}

\textbf{（五）化学分析法（Chemical Analysis）}

收集受试者一日全部食物，在实验室测定营养素含量。样品收集方法：
\begin{itemize}[leftmargin=*]
    \item 双份法：是最准确的方法，制作两份完全相同的食物，一份供食用，另一份作为分析样品。
    \item 收集研究期间消耗的各种未加工食物，从当地市场上购买同样食物作为样品。
\end{itemize}
\begin{itemize}[leftmargin=*]
    \item \textbf{优点：}结果非常准确，能够准确地得出个体营养素准确摄入量。
    \item \textbf{缺点：}过程复杂、成本高，仅适于较小规模调查，如营养代谢试验或某些成分的功能评价研究需要时使用。
\end{itemize}

\textbf{食物量估计工具：}
\begin{enumerate}[leftmargin=*]
    \item 量具（碗、杯、勺等）
    \item 食物图片
    \item 食物模型
    \item 标准餐具
    \item 计算机图像
\end{enumerate}
\end{defbox}

\subsection{体格测量}

\begin{keybox}
\textbf{人体测量常用指标与评价：}

\textbf{1. Broca改良公式（理想体重/标准体重）：}
\[
\text{标准体重(kg)} = \text{身高(cm)} - 105
\]
\[
\text{标准体重} \pm 10\% \text{为正常范围}
\]
\begin{itemize}[leftmargin=*]
    \item 胖瘦 $<$ 10\% 正常
    \item $\pm$ 10\%-20\% 为瘦弱/超重
    \item $>$ 20\% 为严重瘦弱/肥胖
    \item +20\%-30\% 轻度肥胖
    \item +30\%-50\% 中度肥胖
    \item +50\% 重度肥胖
\end{itemize}

\textbf{2. 体质指数（BMI, Body Mass Index）：}
\[
\text{BMI} = \frac{\text{体重(kg)}}{[\text{身高(m)}]^{2}}
\]
我国标准：
\begin{itemize}[leftmargin=*]
    \item BMI $<$ 18.5：消瘦
    \item BMI 18.5-23.9：正常
    \item BMI 24.0-27.9：超重
    \item BMI $\ge$ 28：肥胖
\end{itemize}

\textbf{3. 儿童专用指标：}
\begin{itemize}[leftmargin=*]
    \item 年龄别体重（weight for age）：评价儿童营养状况
    \item 身高别体重（weight for height）：评价儿童营养状况
    \item 年龄别体重（weight for age）：婴幼儿专用
\end{itemize}
\end{keybox}

\begin{defbox}
\textbf{体格测量方法的原理与应用（熟悉）：}
\begin{itemize}[leftmargin=*]
    \item 体格测量用于对大体结构和皮下脂肪的测量，其大小和比例随营养程度的改变而变化。
    \item 对可衡量长期蛋白质和能量摄入不平衡的个体有较高敏感性。学龄儿童的体格测量结果可作为评价群体营养状况的指标。
    \item \textbf{常用指标：}身高（身长）、体重、上臂围、腰围、皮褶厚度。
    \item \textbf{专门指标：}坐高、头围、体成分、骨密度。
    \item 可借助设备和活动辅助测量。
\end{itemize}
\end{defbox}

\subsection{营养水平生化检验与临床检查}

\begin{tipbox}
\textbf{营养水平生化检验（实验室检查方法）（了解）：}
用于评价个体营养素缺乏的风险水平。
\begin{itemize}[leftmargin=*]
    \item 当机体组织中营养素逐渐耗竭，与营养素有关的酶水平和活性或代谢产物开始改变。
    \item 通过测定实验室/生化某些特定营养素的变化，血液学功能改变先于临床症状——出现营养素缺乏。
    \item 常用标本和测定指标：血清（血浆）——如血清铁蛋白、血浆视黄醇、血清25-(OH)D、血清钙等。
\end{itemize}

\textbf{常见营养水平鉴定生化指标（了解）：}
\begin{enumerate}[leftmargin=*]
    \item \textbf{蛋白质：}血浆总蛋白、白蛋白、前白蛋白、运铁蛋白、视黄醇结合蛋白——白蛋白半衰期较长（约20天），前白蛋白和运铁蛋白半衰期短（约2-8天），反映近期蛋白质营养状况更敏感。
    \item \textbf{铁：}血清铁蛋白（反映贮存铁）、血清铁、总铁结合力、运铁蛋白饱和度、血红蛋白。
    \item \textbf{维生素A：}血浆视黄醇水平。
    \item \textbf{维生素D：}血清25-(OH)D水平。
    \item \textbf{维生素B_{1}：}红细胞转酮醇酶活性系数（ETK-AC）。
    \item \textbf{维生素B_{2}：}红细胞谷胱甘肽还原酶活性系数（EGR-AC）。
    \item \textbf{维生素C：}血浆维生素C浓度、白细胞维生素C含量。
    \item \textbf{碘：}尿碘中位数。
    \item \textbf{钙：}血清钙（受甲状旁腺激素精密调控，不敏感）、骨密度测定。
    \item \textbf{锌：}血清/血浆锌浓度。
\end{enumerate}

\textbf{临床检查方法（熟悉）：}
\begin{enumerate}[leftmargin=*]
    \item \textbf{病史采集：}一般情况、目前健康状况、病史、既往病史（吸烟、膳食情况、饮酒和体力活动等、药物、日常体力活动）。
    \item \textbf{与营养素缺乏有关的检查：}各系统——皮肤、肌肉、骨骼、心血管、消化和神经系统的检查。检查部位：头、面部、眼、唇、舌、牙齿和齿龈、指甲。
\end{enumerate}

\textbf{常见营养缺乏症临床表现（了解）：}

\begin{table}[H]
\centering
\begin{tabularx}{\textwidth}{lX}
\toprule
\textbf{缺乏营养素} & \textbf{主要临床表现} \\
\midrule
蛋白质-能量（PEM） & 消瘦（marasmus）、水肿型（kwashiorkor）——生长迟缓、皮下脂肪减少、肌肉萎缩 \\
维生素A & 夜盲症、干眼病、毕脱氏斑（Bitot spots）、毛囊角化过度 \\
维生素D/钙 & 儿童佝偻病（颅骨软化、肋串珠、O/X型腿）、成人骨软化症、骨质疏松 \\
维生素B_{1} & 脚气病（干性：末梢神经炎；湿性：心衰水肿；婴儿型：紫绀心衰） \\
维生素B_{2} & 口角炎、唇炎、舌炎、阴囊皮炎、脂溢性皮炎 \\
烟酸 & 癞皮病（皮炎、腹泻、痴呆——"3D"：Dermatitis, Diarrhea, Dementia） \\
维生素C & 坏血病（牙龈出血、毛囊周围出血、瘀斑、伤口愈合不良） \\
叶酸/维生素B_{12} & 巨幼红细胞性贫血（舌炎、面色苍白、乏力） \\
铁 & 缺铁性贫血（面色苍白、疲乏、口角炎、反甲、异食癖） \\
碘 & 甲状腺肿、克汀病（先天性：严重智力障碍+矮小） \\
锌 & 生长迟缓、性发育延迟、食欲减退、异食癖、皮肤损害 \\
\bottomrule
\end{tabularx}
\end{table}
\end{tipbox}

\subsection{膳食营养评价}

\begin{keybox}
\textbf{膳食评价的依据：}
\begin{enumerate}[leftmargin=*]
    \item \textbf{DRIs：}评价膳食能量和营养素摄入量是否满足需要。
    \item \textbf{膳食指南：}规范膳食行为。
    \item \textbf{食物成分表：}将食物成分转换为营养素成分。
    \item \textbf{烹调加工方法合理性：}评价营养素在加工过程中的损失。
\end{enumerate}

\textbf{膳食营养素摄入量评价标准：}
\begin{itemize}[leftmargin=*]
    \item \textbf{摄入充足：}$>$ 90\% 标准
    \item \textbf{摄入不足：}$<$ 80\% 标准（长期摄入不足导致营养不良）
    \item \textbf{严重缺乏：}$<$ 60\% 标准（严重缺乏可致严重健康损害）
    \item 人群分布评估：关注高危人群比例和缺乏严重程度
\end{itemize}
\end{keybox}

\begin{keybox}
\textbf{能量及营养素评价要点（掌握）：}

\textbf{能量：}
\begin{itemize}[leftmargin=*]
    \item 摄入量与EER比较，总量 $\le$ EER $\pm$ 10\%（对成年人）。
    \item 产能营养素供能比：
    \begin{itemize}
        \item 蛋白质：10\%-15\%（儿童、老人12\%-15\%）
        \item 脂肪：20\%-30\%（儿童25\%-35\%）
        \item 碳水化合物：50\%-65\%
    \end{itemize}
    \item 蛋白质供能比（\%）= 蛋白质摄入量(g) × 4 ÷ 能量总摄入量(kcal) × 100
    \item 碳水化合物供能比（\%）= 碳水化合物摄入量(g) × 4 ÷ 能量总摄入量(kcal) × 100
    \item 脂肪供能比（\%）= 脂肪摄入量(g) × 9 ÷ 能量总摄入量(kcal) × 100
\end{itemize}

\textbf{三餐分配（能量比）：}
\begin{itemize}[leftmargin=*]
    \item 早餐:中餐:晚餐 = 3:4:3
    \item 早餐供能比 = 早餐能量摄入量 ÷ 能量总摄入量 × 100
    \item 中餐供能比 = 中餐能量摄入量 ÷ 能量总摄入量 × 100
    \item 晚餐供能比 = 晚餐能量摄入量 ÷ 能量总摄入量 × 100
\end{itemize}

\textbf{蛋白质：}
\begin{itemize}[leftmargin=*]
    \item 总量 $\ge$ 90\% RNI
    \item 优质蛋白质占 1/3-1/2（动物性蛋白+大豆蛋白）
    \item 来源：动物性食物（肉类、奶类、蛋类等）和豆类及其制品
\end{itemize}

\textbf{脂肪：}
\begin{itemize}[leftmargin=*]
    \item 动物性食物和植物性食物提供的脂肪量
    \item 脂肪酸来源：动物性食物提供的脂肪占总脂肪的百分比（E\%）
    \item 脂肪酸构成：SFA:MUFA:PUFA = 1:1:1
    \item 胆固醇和反式脂肪
\end{itemize}

\textbf{碳水化合物：}
\begin{itemize}[leftmargin=*]
    \item 占总能量百分比，单糖、双糖来源
    \item 复杂碳水化合物来源
    \item 膳食纤维来源
    \item 单糖/双糖与复合糖比例，GI和GL
\end{itemize}

\textbf{矿物质：}
\begin{itemize}[leftmargin=*]
    \item 铁总量 $\ge$ 90\% RNI，来源于动物性食物应占1/2
    \item 钙总量 $\ge$ 90\% RNI，来源于乳制品占2/3
\end{itemize}

\textbf{维生素：}
\begin{itemize}[leftmargin=*]
    \item 维生素A——食物来源与营养素评价
    \item 维生素D——晒太阳和食物来源与营养素评价
    \item 维生素B_{1}、B_{2}、叶酸、维生素C
\end{itemize}
\end{keybox}

\begin{tipbox}
\textbf{一般意义的营养缺乏：}
\begin{itemize}[leftmargin=*]
    \item 儿童生长迟缓：身高低于同年龄标准
    \item 儿童消瘦：体重低于同身高儿童
    \item 儿童超重：体重高于同身高儿童
    \item 成人超重：身体脂肪增多，BMI > 25
    \item 微量营养素缺乏症：如维生素A、锌、碘、叶酸处于较低水平
    \item 成人肥胖症：身体脂肪增多，BMI $\ge$ 30
    \item 非传染性疾病：糖尿病、心脏病、某些癌症
\end{itemize}
\end{tipbox}

% [补充自往届笔记]
平衡膳食五要求（食谱编制的核心原则）：
\begin{enumerate}[leftmargin=*]
    \item 提供种类齐全、数量充足、比例适宜的营养素
    \item 提供数量适当的能量和各种营养素
    \item 食物应新鲜卫生
    \item 正确烹调加工
    \item 良好的进餐制度和环境
\end{enumerate}

% =====================================================================
\section{营养监测与食谱编制}
% =====================================================================

\subsection{营养监测}

\begin{defbox}
\textbf{营养监测（Nutrition Surveillance）：}通过纵向持续动态监测目标人群营养状况，掌握营养状况变化趋势，为制定政策和干预措施提供依据。

\textbf{营养监测与营养调查的区别：}
\begin{itemize}[leftmargin=*]
    \item 营养调查：横断面调查，在某时间点评估目标人群营养状况。
    \item 营养监测：纵向动态监测，持续追踪营养状况变化趋势。
\end{itemize}

\textbf{常用监测指标：}
\begin{enumerate}[leftmargin=*]
    \item \textbf{健康指标：}身高、体重、贫血患病率、低出生体重率、5岁以下儿童死亡率等。
    \item \textbf{食物消费指标：}人均食物消费量、膳食能量和营养素摄入量。
    \item \textbf{社会经济指标：}恩格尔系数（Engel's Coefficient）——食物支出占家庭或个人总消费支出的百分比。
    \begin{itemize}[leftmargin=*]
        \item 恩格尔系数 $>$ 60\%：贫困
        \item 50\%-59\%：温饱
        \item 40\%-49\%：小康
        \item 30\%-39\%：富裕
        \item $<$ 30\%：最富裕
    \end{itemize}
\end{enumerate}
\end{defbox}

\subsection{食谱编制}

\begin{defbox}
\textbf{食谱编制的基本原则、依据和方法（熟悉）：}

\textbf{基本原则：}
\begin{enumerate}[leftmargin=*]
    \item \textbf{参照DRIs，保证营养均衡：}满足用膳者能量和各种营养素需要。
    \item \textbf{食物多样，合理搭配：}充分发挥营养素互补作用。
    \item \textbf{考虑用膳者特点和需求：}年龄、性别、生理状况、劳动强度、疾病状况。
    \item \textbf{注意食品安全和加工合理：}科学烹调，减少营养素损失。
    \item \textbf{考虑经济条件和食物供应：}因地制宜，合理选择食物。
\end{enumerate}

\textbf{编制依据：}
\begin{itemize}[leftmargin=*]
    \item DRIs（作为营养素摄入目标）
    \item 膳食指南和膳食宝塔（作为食物种类和数量参考）
    \item 食物成分表（计算营养素实际含量）
    \item 烹调加工方法的影响
\end{itemize}

\textbf{基本方法：}
\begin{enumerate}[leftmargin=*]
    \item 确定用膳者人数、年龄、性别、生理状况和劳动强度。
    \item 计算能量和营养素需要量（参照DRIs）。
    \item 确定主食（谷类）、副食（蔬菜、肉、蛋、奶、豆等）的种类和数量（参照膳食宝塔）。
    \item 分配一日三餐（早餐30\%、中餐40\%、晚餐30\%能量）。
    \item 编制一日食谱，计算营养素含量，与DRIs比较进行调整。
    \item 编制周食谱，避免单调重复，力求食物多样。
\end{enumerate}

\textbf{食谱编制的工作程序（"日需食谱周排"）：}% [补充自往届笔记]
建议在编制好一日食谱基础上进一步编制周食谱，避免单调重复，力求食物多样。
\end{defbox}

% =====================================================================
\section{公共营养改善}
% =====================================================================

\begin{defbox}
\textbf{公共营养改善（Nutrition Improvements）的定义：}
指采用膳食营养干预措施，改善营养不良人群的营养水平，以促使其恢复营养健康状况的活动。

\textbf{公共营养改善工作：}指为改善居民营养状况而开展的预防和控制营养缺乏、营养过剩和营养相关疾病等工作。

\textbf{公共营养改善常用方法：}
\begin{enumerate}[leftmargin=*]
    \item 营养教育与咨询
    \item 营养配餐与食谱编制
    \item 食品营养标签
    \item 食品营养强化（食品原料、新资源食品等）
\end{enumerate}

\textbf{营养干预（Nutrition Intervention）：}健康促进活动，目的是改变目标人群营养状况，包括选择干预策略、研究目标人群、研究消费者行为和实施评价研究。
\end{defbox}

\subsection{营养教育}

\begin{defbox}
\textbf{营养教育（Nutrition Education）的定义：}
通过改变人们的饮食行为达到改善营养目的的一种有计划的活动。意义：多途径、低成本、覆盖面广；有计划、有组织、有系统。

\textbf{营养教育的作用：}
\begin{enumerate}[leftmargin=*]
    \item 提供膳食行为改变所需的知识、技能和服务。
    \item 帮助个体和群体树立食品与营养的健康意识。
    \item 培养良好膳食行为与生活方式，使人们在面临营养与食品安全问题时有能力做出有益于健康的选择。
\end{enumerate}

\textbf{营养教育的主要内容：}% [补充自往届笔记]
\begin{enumerate}[leftmargin=*]
    \item 营养基础知识
    \item 健康生活方式
    \item 《中国居民膳食指南》和《中国居民平衡膳食宝塔》
    \item 我国人群膳食营养与健康状况及变化趋势
    \item 膳食营养相关的慢性疾病的膳食预防措施
    \item 营养相关的法律、法规与政策
\end{enumerate}

\textbf{营养教育的主要实施方式（了解）：}% [补充自往届笔记]
\begin{enumerate}[leftmargin=*]
    \item 各类人员知识培训：有计划地对餐饮业、农业、商业、轻工、医疗卫生、疾病控制等部门有关人员进行营养知识培训。
    \item 营养科学纳入中小学教育：将营养知识纳入中小学的教育内容和教学计划，安排一定课时的营养知识教育，使学生懂得平衡膳食的原则，培养学生良好的饮食习惯，提高自我保护意识。
    \item 社区教育网络体系：提高社区基层卫生人员和居民的膳食营养意识，培养良好的饮食行为，改善营养状况。
    \item 全社会宣传教育：利用多种宣传媒介，广泛开展群众性营养宣传活动，倡导健康膳食模式和健康生活方式，纠正不良饮食习惯等。
\end{enumerate}
\end{defbox}

\subsection{营养信息传播与营养行为干预}

\begin{tipbox}
\textbf{营养信息传播（了解）：}
\begin{itemize}[leftmargin=*]
    \item 准确信息是改变行为的基石。
    \item 营养信息传播是健康教育的一个重要组成部分，是营养教育和营养改善行动的重要手段和策略。
\end{itemize}

\textbf{营养行为干预（熟悉）：}
\begin{itemize}[leftmargin=*]
    \item 行为改变是实现营养教育计划目标的重要手段。
    \item 通过行为指导、技能训练和辅导咨询，使受教育者实现特定的饮食行为的改变。
    \item 主要方法：模拟、示范、实际操作、个别指导、小组讨论等均属干预手段。此外，还包括一些行为矫正技术。
\end{itemize}
\textbf{营养干预应用（了解）：}纳入初级卫生保健服务体系，利用多种宣传媒介开展科普宣传，将营养知识纳入中小学教育。
\end{tipbox}

\subsection{食品营养强化}

\begin{defbox}
\textbf{食品营养强化 / 食品强化（Food Fortification）：}
根据不同人群营养需要，向食品中添加一种或多种营养素（或某些天然食物成分），以提高食品营养价值的过程称为食品营养强化，或简称食品强化。

\textbf{食品强化剂：}指为增强营养成分而加入食品中的天然或者人工合成的属于天然营养素范围的食品添加剂。主要包括必需氨基酸、维生素、矿物质等。

\textbf{食品营养强化的意义：}
\begin{enumerate}[leftmargin=*]
    \item 弥补天然食物的营养缺陷，补充其不足（如谷类加赖氨酸）。
    \item 补充加工过程中人为损失的营养素（B族维生素、精白米面等）。
    \item 满足特定人群的营养需要（配方奶粉、要素膳等）。
    \item 预防营养不良（食盐加碘）。
    \item 改善食物感官性状和营养价值：婴儿配方食品。
\end{enumerate}
\end{defbox}

% [补充自往届笔记]
我国食物营养相关政策及规划（了解）：
\begin{itemize}[leftmargin=*]
    \item 食品安全法
    \item 营养改善工作管理办法
    \item 全国营养调查工作规范
    \item 食品营养标签管理规范
    \item 《2025-2030年中国食物与营养发展纲要》
    \item 2017年7月《国民营养计划（2017-2030）》
\end{itemize}

% =====================================================================
\section{食品营养标签、NRV与功能食品}
% =====================================================================

\subsection{食品营养标签}

\begin{defbox}
\textbf{食品营养标签（Food Nutrition Label）（熟悉）：}
食品营养标签是向消费者提供食品营养信息和特性的说明，包括营养成分表、营养声称和营养成分功能声称。食品营养标签是预包装食品标签的一部分，管理遵循相关法规（如《GB 28050-2011 食品安全国家标准 预包装食品营养标签通则》）。% [补充自往届笔记] 强制标示内容包括：能量、蛋白质、脂肪、碳水化合物和钠（"1+4"核心营养素）。

\textbf{营养成分表基本内容：}
\begin{itemize}[leftmargin=*]
    \item 能量和核心营养素（1+4）：能量 + 蛋白质、脂肪、碳水化合物、钠（强制标示）。
    \item 营养声称（如"高钙"、"低脂"、"富含膳食纤维"等）。
    \item 营养成分功能声称（如"钙有助于骨骼和牙齿的健康"）。
\end{itemize}
\end{defbox}

\subsection{NRV（营养素参考值）}

\begin{keybox}
\textbf{营养素参考值（NRV——Nutrient Reference Values）（掌握）：}
\begin{itemize}[leftmargin=*]
    \item \textbf{定义：}专用于食品营养标签上比较食品营养素含量的参考值。
    \item \textbf{用途：}NRV是食品营养标签上的"标尺"，用NRV\%表示食品中某营养素含量占NRV的百分比，帮助消费者快速判断该食品在每日膳食中的贡献比例，便于比较不同食品营养价值，做出合理食物选择。
    \item \textbf{示例：}某食品每份含钙200 mg，钙NRV为800 mg，则钙NRV\% = 200/800 $\times$ 100\% = 25\%。
    \item \textbf{NRV与DRIs的区别：}DRIs是一组针对不同年龄、性别、生理状况人群的参考值体系；NRV是用于营养标签的单一参考值，通常采用成人RNI或AI。NRV用于食品标签，DRIs用于膳食评价和规划。
\end{itemize}
\end{keybox}

\subsection{功能食品}

\begin{defbox}
\textbf{功能食品（Functional Food）的概念（熟悉）、分类（了解）、功能评价（了解）及管理（熟悉）：}

\textbf{概念：}功能食品是指具有特定营养保健功能的食品，即适宜于特定人群食用，具有调节机体功能，不以治疗疾病为目的的食品。在我国与"保健食品"概念基本对应。

\textbf{分类（了解）：}
\begin{enumerate}[leftmargin=*]
    \item 营养素补充剂——补充维生素、矿物质等。
    \item 具有特定保健功能的食品——按功能声称分类（如增强免疫力、辅助降血脂、辅助降血糖、抗氧化、辅助改善记忆、缓解视疲劳、促进排铅、清咽、辅助降血压、改善睡眠、促进泌乳、缓解体力疲劳、提高缺氧耐受力、对辐射危害有辅助保护功能、减肥、改善生长发育、增加骨密度、改善营养性贫血、对化学性肝损伤有辅助保护功能、祛痤疮、祛黄褐斑、改善皮肤水分、改善皮肤油分、调节肠道菌群、促进消化、通便、对胃黏膜有辅助保护功能等27类保健功能）。
\end{enumerate}

\textbf{功能评价（了解）：}功能食品需进行功能学实验验证和安全性评价。功能学实验包括动物功能实验和/或人体试食试验，证明其具有声称的保健功能。安全性毒理学评价根据原料不同进行相应的毒理学试验。

\textbf{管理（熟悉）：}我国对保健食品实行注册与备案双轨制管理。需经国家市场监督管理总局审批，获得"蓝帽子"标志（保健食品标志）后方可上市销售。产品标签和说明书必须标注保健功能、适宜人群、不适宜人群、食用方法和用量等，须标明"本品不能替代药物"。
\end{defbox}

% =====================================================================
\reviewcontent{
\section{复习题}
% =====================================================================

\begin{enumerate}[leftmargin=*, itemsep=8pt]
    \item \textbf{【名词解释】}DRIs；EAR；RNI；AI；UL；AMDR；PI-NCD；SPL；BMI；NRV；恩格尔系数；EER；公共营养；营养调查；营养监测；营养教育
    \item \textbf{【简答题】}简述公共营养的定义、地位与工作内容。
    \item \textbf{【简答题】}简述DRIs七项指标体系的定义、制定目标和相互关系。
    \item \textbf{【简答题】}比较AI与RNI的异同以及UL的制定注意事项。
    \item \textbf{【简答题】}简述DRIs的制定修订程序及其用途层次。
    \item \textbf{【简答题】}列出《中国居民膳食指南(2022)》8条核心准则及膳食指南对群体的核心要求。
    \item \textbf{【简答题】}简述膳食结构的四种类型、各自特点及其与健康的关系，以及我国膳食结构现状。
    \item \textbf{【简答题】}简述营养调查的目的、内容和营养状况综合评价的四个方面。
    \item \textbf{【简答题】}比较膳食调查五种方法（称重法、24h回顾法、FFQ、化学分析法、记账法）的特点、优缺点及应用。
    \item \textbf{【简答题】}简述膳食调查方法选择依据（5W1H原则）和基本要求。
    \item \textbf{【简答题】}简述膳食营养评价的依据、摄入量评价标准和主要营养素评价要点。
    \item \textbf{【简答题】}简述BMI分级标准、Broca改良公式及其应用。
    \item \textbf{【简答题】}简述营养监测与营养调查的概念区别，以及恩格尔系数的含义和分级标准。
    \item \textbf{【简答题】}简述食谱编制的基本原则、依据和方法。
    \item \textbf{【简答题】}简述公共营养改善常用方法、营养教育的概念和主要内容、营养干预的主要方法。
    \item \textbf{【简答题】}简述食品营养强化的概念、食品强化剂的定义及营养强化的意义。
    \item \textbf{【简答题】}简述食品营养标签的基本内容和NRV的概念、用途及其与DRIs的区别。
    \item \textbf{【简答题】}简述功能食品（保健食品）的概念、分类、功能评价及管理要求。
\end{enumerate}

\subsection*{参考答案要点}

\begin{enumerate}[leftmargin=*, itemsep=5pt]
    \item 见正文相关定义。
    \begin{itemize}
        \item \textbf{公共营养：}从医学状况+人群社会视角的群体概念。通过膳食提供给机体种类齐全、数量充足、比例适宜的能量和各种营养素，并与机体需要相平衡。
        \item \textbf{DRIs：}一组每日平均膳食营养素摄入量的参考值，用于评价膳食营养素供给量能否满足人体需要及是否存在过量摄入风险。
        \item \textbf{EAR：}满足群体中50\%个体需要量，是制订RNI的基础。
        \item \textbf{RNI：}满足97\%-98\%个体需要量，RNI = EAR + 2SD 或 RNI = 1.2 × EAR。
        \item \textbf{AI：}资料不足无法计算EAR时，通过观察或实验获得。AI主要用于个体营养素摄入目标，不能用于人群营养评价；准确度不如RNI。
        \item \textbf{UL：}平均每日摄入营养素的最高限量。摄入量 < UL对健康人群无副作用，> UL不良反应风险增加。
        \item \textbf{AMDR：}宏量营养素可接受范围——脂肪20\%-30\%，碳水化合物50\%-65\%，蛋白质10\%-15\%。
        \item \textbf{PI-NCD：}预防非传染性慢性病的建议摄入量，以NCD一级预防为目标。PI可能高于RNI/AI（如维生素C、钾）或低于AI（如钠）。
        \item \textbf{SPL：}特定建议值，专用于预防NCD，适用于其他食物成分（如植物化学物）的预防NCD建议量。
        \item \textbf{EER：}能长期保持良好健康状态、维持体重和机体构成以及体力活动水平的个体及人群，达到能量平衡时所需要的膳食能量摄入量。
        \item \textbf{BMI：}体质指数，体重(kg) / [身高(m)]^{2}。标准：<18.5消瘦，18.5-23.9正常，24.0-27.9超重，≥28肥胖。
        \item \textbf{NRV：}营养素参考值，专用于食品营养标签上比较食品营养素含量的参考值。
        \item \textbf{恩格尔系数：}食物支出占总消费支出百分比。>60\%贫困，50\%-59\%温饱，40\%-49\%小康，30\%-39\%富裕，<30\%最富裕。
        \item \textbf{营养调查：}运用各种手段准确了解某一人群特定时间的各种营养指标水平，判断其营养和健康状况。
        \item \textbf{营养监测：}通过纵向持续动态监测目标人群营养状况，掌握变化趋势。
        \item \textbf{营养教育：}通过改变人们的饮食行为达到改善营养目的的一种有计划的活动。
    \end{itemize}

    \item 公共营养定义、地位与工作内容：
    \begin{itemize}
        \item \textbf{定义：}从医学状况+人群社会视角的群体概念。合理营养（平衡膳食）：提供种类齐全、数量充足、比例适宜的能量和各种营养素。
        \item \textbf{地位：}第一，促进全人群营养与健康；第二，增强国民体质和综合国力；第三，事关国家发展战略问题。
        \item \textbf{工作内容：}营养标准制定（DRIs）、膳食指南制定、人员培训、营养调查与评价、营养监测、营养干预、食物营养规划与立法、其他营养公众指导。
        \item \textbf{营养不良：}营养缺乏（nutrition deficiency）和营养过剩（nutrition excess）。
    \end{itemize}

    \item DRIs七项指标体系：
    \begin{itemize}
        \item \textbf{EAR（平均需要量）：}满足50\%个体需要，是制订RNI的基础。目标：预防缺乏。
        \item \textbf{RNI（推荐摄入量）：}满足97\%-98\%个体需要，RNI = EAR + 2SD 或 1.2 × EAR。目标：预防缺乏。
        \item \textbf{AI（适宜摄入量）：}EAR无法计算时使用，通过观察或实验获得，准确度不如RNI。目标：预防缺乏。
        \item \textbf{UL（可耐受最高摄入量）：}安全上限。目标：防止过量。
        \item \textbf{AMDR（宏量营养素可接受范围）：}脂肪20\%-30\%，碳水化合物50\%-65\%，蛋白质10\%-15\%。目标：预防慢性病。
        \item \textbf{PI-NCD（预防慢病建议摄入量）：}以NCD一级预防为目标，与RNI/AI比较：PI>RNI/AI时以PI为准（维生素C、钾），PI<AI时以AI为准（钠）。目标：预防慢性病。
        \item \textbf{SPL（特定建议值）：}针对其他膳食成分（如植物化学物）的预防NCD建议量。目标：预防慢性病。
    \end{itemize}

    \item AI与RNI异同：
    \begin{itemize}
        \item \textbf{相同点：}都可以作为群体中个体营养素摄入量的目标值，满足群体中几乎所有个体的需要。
        \item \textbf{不同点：}AI准确度不如RNI，主要用于个体营养素摄入目标；群体摄入量≥AI时，摄入不足者危险性很小。
    \end{itemize}
    UL制定注意事项：随强化食品和使用补充剂的发展需制定ULs；毒副作用与摄入总量有关的，按食物+饮水+补充剂总量制定；毒副作用仅与强化食物和补充剂有关的，按这些来源制定；未定UL不意味着过多摄入无潜在危害。

    \item DRIs制定修订程序：人群调查监测→制定DRIs→膳食指导→食物指导→食品营养标签→人群营养改善反馈→营养科学研究→修订DRIs。每4-5年修订一次。用途层次：个体（膳食规划、食谱、配方研发）、群体（膳食评价、营养咨询）、安全作用（标准/依据）。

    \item 膳食指南8条核心准则（2022）：
    \begin{enumerate}[leftmargin=*]
        \item 食物多样，合理搭配（每天>12种，每周>25种）
        \item 吃动平衡，健康体重
        \item 多吃蔬果、奶类、全谷、大豆（奶300-500 ml/d）
        \item 适量吃鱼、禽、蛋、瘦肉（120-200 g/d，每周鱼2次）
        \item 少盐少油，控糖限酒（盐<5 g/d，糖<50 g/d最好<25 g/d）
        \item 规律进餐，足量饮水（女1500 ml，男1700 ml）
        \item 会烹会选，会看标签
        \item 公筷分餐，杜绝浪费
    \end{enumerate}
    膳食宝塔五层：第一层谷薯类250-400 g/d；第二层蔬菜300-500 g/d、水果200-350 g/d；第三层鱼禽肉蛋120-200 g/d；第四层奶及奶制品300-500 g/d、大豆及坚果25-35 g/d；第五层盐<5 g/d、油25-30 g/d。饮水1500-1700 ml/d，活动6000步/d。
    膳食指南对群体的核心要求：群体性、科学性、普及性、有效性、通俗性、权威性、保障性、预防性。

    \item 四种膳食模式：
    \begin{itemize}
        \item \textbf{动物性食物为主型（西方模式）：}"三高一低"（高能量、高脂肪、高蛋白、低膳食纤维），与肥胖、心脑血管疾病、结肠癌直肠癌等慢性病风险增加相关。
        \item \textbf{植物性食物为主型（东方模式）：}"一高三低"（低蛋白质、低脂肪、植物性食物为主），感染性疾病+慢性病并存。
        \item \textbf{动植物平衡型（日本模式）：}动植物消费均衡，水产品消费量大，三大宏量营养素供能比合理。
        \item \textbf{地中海模式：}橄榄油为主，不饱和脂肪高，大量蔬菜水果全谷物豆类，适量鱼类和红酒。
    \end{itemize}
    我国膳食结构：营养缺乏与过剩并存，地域偏差明显，正朝富裕型膳食结构转变。

    \item 营养调查目的6项（见正文）；内容4项：膳食调查、体格测量、生化检验、临床检查。

    \item 五种方法比较：
    \begin{itemize}
        \item \textbf{称重法：}准确可靠，3-7天，可作为"金标准"，费时费力，不适于大规模调查。
        \item \textbf{24h回顾法：}最常用，不改变饮食习惯，简便快速，连续3天（2工作日+1周末），依赖记忆，对某些人群不适用。
        \item \textbf{FFQ：}可调查长期（几天到1年以上）膳食模式，方法简单费用低，精确度较低，分定性/定量/半定量三种。与24h回顾法联用更全面。
        \item \textbf{化学分析法：}最准确的方法（双份法），过程复杂成本高，仅适用于小规模。
        \item \textbf{记账法：}可调查较长时期（一个月或更长），费用低，不能分析个体摄入状况。
    \end{itemize}

    \item 膳食调查方法选择依据（5W1H）：Who（个体/群体）、What（食物/营养素）、When（当前/通常模式/长期）、Where（在家/在外）、Why（研究目的），以及研究人群大小、精确度要求、经费及时间。基本要求：观察单位、资料收集方式、研究时限、食物量估计、食物-营养素转换方法。

    \item 膳食营养评价依据：DRIs（摄入充足>90\%标准、摄入不足<80\%标准、严重缺乏<60\%标准）、膳食指南、食物成分表、烹调加工方法合理性。评价要点：能量EER±10\%；蛋白质总量≥90\%RNI优质蛋白1/3-1/2供能比10\%-15\%；脂肪供能比20\%-30\%、SFA:MUFA:PUFA=1:1:1；碳水化合物50\%-65\%；三餐分配3:4:3；铁≥90\%RNI动物来源1/2；钙≥90\%RNI乳制品来源2/3。

    \item BMI分级标准：<18.5消瘦，18.5-23.9正常，24.0-27.9超重，≥28肥胖。Broca改良公式：标准体重(kg)=身高(cm)-105。±10\%正常，±10\%-20\%瘦弱/超重，>20\%严重瘦弱/肥胖，+20\%-30\%轻度肥胖，+30\%-50\%中度肥胖，+50\%重度肥胖。

    \item 营养调查为横断面调查，营养监测为纵向动态监测。恩格尔系数：食物支出/总消费支出×100\%。>60\%贫困，50\%-59\%温饱，40\%-49\%小康，30\%-39\%富裕，<30\%最富裕。

    \item 食谱编制基本原则：满足DRIs、食物多样、考虑用膳者特点和需求、注意食品安全科学加工、考虑经济条件和食物供应。依据：DRIs、膳食指南和宝塔、食物成分表、烹调加工方法。方法：确定人群特点→计算营养素需要量→确定主副食种类数量→分配三餐（3:4:3）→编制食谱计算营养素→与DRIs比较调整→编制周食谱。

    \item 公共营养改善常用方法：营养教育与咨询、营养配餐与食谱编制、食品营养标签、食品营养强化。营养教育：通过改变饮食行为改善营养的有计划活动（多途径、低成本、覆盖面广）。主要内容：营养基础知识、健康生活方式、膳食指南和宝塔、人群营养与健康状况、慢性病膳食预防、营养法律法规。营养干预方法：模拟、示范、实际操作、个别指导、小组讨论、行为矫正技术等。

    \item 食品营养强化：向食品中添加营养素以提高食品营养价值的过程。食品强化剂：天然或人工合成的属于天然营养素范围的食品添加剂（必需氨基酸、维生素、矿物质等）。意义5项：弥补天然食物营养缺陷；补充加工损失；满足特定人群需要；预防营养不良；改善食物感官性状和营养价值。

    \item NRV：专用于食品营养标签的参考值。NRV\%表示营养素含量占NRV的百分比，帮助消费者判断该食品在每日膳食中的贡献比例。NRV与DRIs区别：DRIs是针对不同人群的参考值体系，NRV是用于标签的单一参考值（取成人RNI/AI）。营养成分表强制标示"1+4"：能量+蛋白质、脂肪、碳水化合物、钠。

    \item 功能食品（保健食品）：具有特定保健功能、适宜特定人群食用、不以治疗疾病为目的的食品。分类：营养素补充剂和27类特定保健功能食品。管理：注册与备案双轨制，国家市场监督管理总局审批，获"蓝帽子"标志，标签须标明保健功能、适宜人群、不适宜人群、食用方法，须标明"本品不能替代药物"。功能评价：功能学实验（动物实验和/或人体试食试验）和安全性毒理学评价。
\end{enumerate}

}
\newpage
